\documentclass[10pt]{article}

% ---------- Document Identity (update these only) ----------
\newcommand{\DocStage}{}
\newcommand{\DocTitle}{Your Road to Tier One SOF}
\newcommand{\ProgramName}{182-Week Civilian-to-Selection Readiness Pipeline}
\newcommand{\ProgramSubtitle}{}

% ---------- Page ----------
\usepackage[legalpaper,left=0.5in,right=0.5in,top=0.6in,bottom=0.45in]{geometry}

% ---------- Typography ----------
\usepackage[T1]{fontenc}
\usepackage[utf8]{inputenc}
\usepackage{libertinus}
\usepackage{microtype}
\usepackage{parskip}
\usepackage{ragged2e}
\usepackage{textcomp}
\AtBeginDocument{\color{black}}
\AtBeginDocument{\pagecolor{white}}
\AtBeginDocument{\justifying}

% ---------- Landscape pages ----------
\usepackage{pdflscape}

% ---------- Color ----------
\usepackage[table]{xcolor}
\definecolor{StageBlue}{HTML}{1F4E79}
\definecolor{StageBlueDark}{HTML}{163A5A}
\definecolor{LightGray}{HTML}{F2F4F7}
\definecolor{MidGray}{HTML}{D9DEE5}
\definecolor{RuleGray}{HTML}{C7CED8}
\definecolor{HeaderInk}{HTML}{000000}
\definecolor{Ink}{HTML}{000000}

% ---------- Utility ----------
\usepackage{enumitem}
\sloppy
\setlength{\emergencystretch}{2em}

% ---------- Boxes / UI ----------
\usepackage[most]{tcolorbox}

\tcbset{
  charterTitle/.style={
    enhanced,
    colback=StageBlue!4,
    colframe=StageBlue,
    boxrule=1.25pt,
    arc=3pt,
    left=16pt,right=16pt,top=14pt,bottom=14pt,
    borderline north={3.2pt}{0pt}{StageBlue},
    borderline west={4pt}{0pt}{StageBlue},
    borderline south={1.25pt}{0pt}{StageBlue},
    drop shadow={black!25!white},
  },
  charterRule/.style={
    enhanced,
    colback=LightGray,
    colframe=RuleGray,
    boxrule=0.85pt,
    arc=3pt,
    left=12pt,right=12pt,top=10pt,bottom=10pt,
    borderline west={3pt}{0pt}{StageBlue},
  },
  charterBadge/.style={
    enhanced,
    colback=StageBlue,
    coltext=white,
    colframe=StageBlueDark,
    boxrule=0.6pt,
    arc=2pt,
    left=6pt,right=6pt,top=3pt,bottom=3pt,
  }
}
\newtcbox{\Badge}{nobeforeafter,charterBadge}

% ---------- Lists ----------
\setlist[itemize]{leftmargin=1.5em, itemsep=0.25em, topsep=0.25em}
\setlist[enumerate]{leftmargin=1.8em, itemsep=0.25em, topsep=0.25em}

% ---------- TOC ----------
\usepackage{tocloft}
\setcounter{tocdepth}{3}
\setlength{\cftbeforesecskip}{3pt}
\setlength{\cftbeforesubsecskip}{2pt}
\setlength{\cftbeforesubsubsecskip}{0pt}
\renewcommand{\cftsecfont}{\bfseries}
\renewcommand{\cftsecpagefont}{\bfseries}
\renewcommand{\cftsubsecfont}{\normalfont}
\renewcommand{\cftsubsecpagefont}{\normalfont}
\renewcommand{\cftsubsubsecfont}{\normalfont}
\renewcommand{\cftsubsubsecpagefont}{\normalfont}
\setlength{\cftsecindent}{0pt}
\setlength{\cftsubsecindent}{1.25em}
\setlength{\cftsubsubsecindent}{2.8em}
\setlength{\cftsecnumwidth}{2.1em}
\setlength{\cftsubsecnumwidth}{2.8em}
\setlength{\cftsubsubsecnumwidth}{3.6em}

% Remove the built-in TOC heading so only the custom heading is shown
\makeatletter
\renewcommand{\tableofcontents}{\@starttoc{toc}}
\makeatother

% ---------- Header/Footer: page number only ----------
\usepackage{fancyhdr}
\pagestyle{fancy}
\fancyhf{}
\renewcommand{\headrulewidth}{0pt}
\renewcommand{\footrulewidth}{0pt}
\fancyfoot[C]{\thepage}

% ---------- Section formatting ----------
\usepackage{titlesec}

\titleformat{\section}
  {\Large\bfseries\color{Ink}}
  {\thesection}{0.6em}{}
  [\vspace{0.25em}\textcolor{StageBlue}{\rule{\linewidth}{0.8pt}}]

\titleformat{\subsection}
  {\large\bfseries\color{Ink}}
  {\thesubsection}{0.6em}{}

\titleformat{\subsubsection}
  {\normalsize\bfseries\color{Ink}}
  {\thesubsubsection}{0.6em}{}

\titlespacing*{\section}{0pt}{0.95em}{0.35em}
\titlespacing*{\subsection}{0pt}{0.65em}{0.25em}
\titlespacing*{\subsubsection}{0pt}{0.55em}{0.2em}

% ---------- Table engine ----------
\usepackage{tabularray}
\UseTblrLibrary{booktabs}

% ---------- tabularray: GLOBAL table treatment ----------
\SetTblrInner{
  cells = {fg=black, bg=white, valign=t},
}

% ---------- Header helpers ----------
\newcommand{\Hdr}[1]{\shortstack[c]{\strut #1\strut}}

% ---------- Lines ----------
\newcommand{\TitleLine}{\textcolor{StageBlue}{\rule{0.98\linewidth}{0.95pt}}}
\newcommand{\SubTitleLine}{\textcolor{RuleGray}{\rule{0.98\linewidth}{0.65pt}}}

% ---------- Continued on next page ----------
\newcommand{\ContinuedOnNextPageInline}{%
  \par
  \vfill
  \noindent\textcolor{RuleGray}{\rule{\linewidth}{1pt}}\vspace{-2pt}
  \noindent\hfill\textit{Continued on next page}\par\vspace{-10pt}
  \noindent\textcolor{RuleGray}{\rule{\linewidth}{1pt}}
  \clearpage
}

% ---------- PDF hyperlinks ----------
\usepackage[hidelinks]{hyperref}
\hypersetup{
  colorlinks=false,
  pdfborder={0 0 0},
  linktoc=all,
  pdftitle={\DocTitle},
  pdfauthor={\ProgramName}
}

% =========================================================
\begin{document}

\section*{7.05 — Cognitive Performance and Energy Enhancers — OTC}
\phantomsection
\addcontentsline{toc}{section}{7.05 — Cognitive Performance and Energy Enhancers — OTC}

\subsection*{7.05.1 Purpose \& Scope}
\phantomsection
\addcontentsline{toc}{subsection}{7.05.1 Purpose \& Scope}

7.05 defines optional OTC cognitive/energy enhancer capability posture and the governance rules that constrain it. The intent is to reduce preventable execution failures (fatigue-driven cancellations, inconsistent alertness) without permitting stimulant misuse, sleep displacement, anxiety escalation, or evidence contamination in protected windows. This overview frames stimulants as a controlled variable, not a performance lever and not a substitute for recovery fundamentals.

\subsection*{7.05.2 Governance and Safety Boundaries}
\phantomsection
\addcontentsline{toc}{subsection}{7.05.2 Governance and Safety Boundaries}

\begin{itemize}
\item OTC-only posture. No prescription guidance of any kind.
\item Safety hierarchy is non-negotiable: sleep → nutrition/hydration → training dose control → clinician evaluation override any energy product considerations.
\item Prohibited product classes (non-negotiable):
  \begin{itemize}
  \item proprietary stimulant blends
  \item fat burners / “thermogenics”
  \item DMAA/DMHA-like stimulants
  \item ephedra-like products
  \item undisclosed stimulant products
  \item “research chemical” or gray-market nootropics
  \end{itemize}
\item Change-control posture:
  \begin{itemize}
  \item one item at a time (introduce or remove only one cognitive/energy variable per change window),
  \item no changes near Deload+Test weeks or gate windows,
  \item log start/stop and adverse effects; treat unlogged changes as validity-contaminating context.
  \end{itemize}
\item Prescription medication interaction caution posture (referral-only, non-diagnostic):
  \begin{itemize}
  \item if already on prescription stimulants or blood pressure medication, treat OTC stimulant use as clinician-led only (do not self-experiment; do not combine; discuss with Medical).
  \end{itemize}
\item Validity posture:
  \begin{itemize}
  \item stimulants do not upgrade readiness; \textbf{VALID} evidence still controls gates.
  \item do not use novel stimulants to “salvage” test performance or to justify advancement claims.
  \end{itemize}
\end{itemize}

\subsection*{7.05.3 Tier Doctrine Applied to Cognitive/Energy Enhancers}
\phantomsection
\addcontentsline{toc}{subsection}{7.05.3 Tier Doctrine Applied to Cognitive/Energy Enhancers}

\begin{itemize}
\item Tier 0: household caffeine sources and non-product fundamentals. Tier 0 is “do nothing new” posture and is the default for most of the pipeline.
\item Tier 1: single-ingredient, transparent-format OTC options that can be logged and controlled with minimal interaction complexity.
\item Tier 2: optional expanded cognitive support options emphasizing transparency, third-party testing, and conservative risk posture (still no stacks).
\item Tier 3: overbuilt governance and documentation posture (verification preference rules, standardization packs), not “stronger stimulants.”
\end{itemize}

Tiering in 7.05 describes governance complexity and controllability, not “required” capability. A higher tier is never framed as necessary for phase completion.

\subsection*{7.05.4 Selection Criteria \& Fit Checks}
\phantomsection
\addcontentsline{toc}{subsection}{7.05.4 Selection Criteria \& Fit Checks}

\begin{itemize}
\item Selection criteria (generic; non-prescriptive):
  \begin{itemize}
  \item single-ingredient formats are preferred over blends to reduce interaction uncertainty and to support one-change-at-a-time posture.
  \item avoid products with “proprietary matrix,” “energy blend,” or undisclosed stimulant labeling.
  \item prefer third-party testing/COA posture when available; keep label transparency as a minimum.
  \end{itemize}
\item Fit checks (operational):
  \begin{itemize}
  \item sleep stability preserved (no new insomnia or sleep fragmentation trend).
  \item no anxiety/panic escalation, agitation, or irritability spikes.
  \item no palpitations, dizziness, lightheadedness, or unusual breathlessness at low effort.
  \item no GI distress that affects session execution.
  \item stable, repeatable response across multiple non-protected weeks before any decision window.
  \end{itemize}
\item “Not fit” triggers:
  \begin{itemize}
  \item panic symptoms or severe anxiety spike.
  \item palpitations with dizziness/lightheadedness.
  \item repeated insomnia or sleep collapse pattern.
  \item any tendency to increase use to compensate for poor sleep or poor recovery.
  \end{itemize}
\end{itemize}

\subsection*{7.05.5 Setup, Safety, and Anchoring}
\phantomsection
\addcontentsline{toc}{subsection}{7.05.5 Setup, Safety, and Anchoring}

\begin{itemize}
\item Anchoring rules:
  \begin{itemize}
  \item stimulants are optional; they do not create readiness and do not justify gate claims.
  \item sleep and training dose control are the primary levers; if fatigue is high, downshift rather than stimulate.
  \end{itemize}
\item Change-control implementation:
  \begin{itemize}
  \item introduce only one cognitive/energy variable at a time.
  \item do not introduce new products or change format/timing near Deload+Test weeks.
  \item maintain stable patterns during decision windows (no “new coffee,” “new gum,” “new tablet,” “new drink” changes).
  \end{itemize}
\item Prescription medication context posture (non-diagnostic; referral-only):
  \begin{itemize}
  \item if on prescription stimulants or blood pressure medication: clinician-led only for any OTC stimulant consideration; default to omit unless cleared.
  \item if on psychiatric medications: treat stimulant/activating products as higher-risk; clinician discussion when uncertain; omit if ambiguous cases.
  \end{itemize}
\end{itemize}

\subsection*{7.05.6 Maintenance, Replacement, and Storage}
\phantomsection
\addcontentsline{toc}{subsection}{7.05.6 Maintenance, Replacement, and Storage}

\begin{itemize}
\item Storage:
  \begin{itemize}
  \item keep products dry and within safe temperature ranges.
  \item avoid leaving consumables in vehicles during extreme temperature swings.
  \end{itemize}
\item Replacement:
  \begin{itemize}
  \item track expiration dates; discard expired products.
  \item maintain like-for-like replacement during decision windows; do not change formulation mid-window.
  \end{itemize}
\item Documentation hygiene:
  \begin{itemize}
  \item keep a stable log record of: product type, format (coffee/tea/tablet/gum), start date, stop date, observed effects, adverse effects.
  \item if lot/batch information is available, keep label photo for audit traceability.
  \end{itemize}
\end{itemize}

\subsection*{7.05.7 Substitution Rules — Fundamentals First Alternatives}
\phantomsection
\addcontentsline{toc}{subsection}{7.05.7 Substitution Rules — Fundamentals First Alternatives}

\begin{itemize}
\item Fundamentals-first substitutions (preferred):
  \begin{itemize}
  \item sleep opportunity protection, schedule adjustments, and downshift posture take precedence over adding stimulants.
  \item hydration and nutrition adequacy should be corrected first; do not misattribute dehydration or under-fueling as “need caffeine.”
  \item training dose control: reduce intensity and/or select lower-impact modalities when fatigue is high; do not stimulate to preserve a planned intensity day.
  \end{itemize}
\item Validity/logging linkage (Stage 2.E posture at high level):
  \begin{itemize}
  \item if a stimulant variable changes during a decision window, it is logged as a controlled variable and gate interpretation becomes variance-sensitive.
  \item if adverse effects cause missed or aborted tests, the test is treated as not admissible for gate credit and is reslotted under the protected-window posture.
  \end{itemize}
\item If a hard stop trigger occurs (panic/anxiety spike, palpitations, dizziness):
  \begin{itemize}
  \item stop the stimulant variable,
  \item log the event,
  \item do not test that day,
  \item apply conservative reslot posture and clinician referral when indicated.
  \end{itemize}
\end{itemize}

\subsection*{7.05.8 Phase Integration Map — Required}
\phantomsection
\addcontentsline{toc}{subsection}{7.05.8 Phase Integration Map — Required}

\begingroup
\scriptsize
\setlength{\tabcolsep}{2pt}
\renewcommand{\arraystretch}{1.10}

\begin{longtblr}{
  width=\linewidth,
  rowhead=1,
  colspec = {
    Q[c,wd=1.05in]
    Q[c,wd=1.55in]
    X[c,2.85]
    X[c,3.85]
  },
  hlines,
  vlines,
  cells = {hsep=2pt, vsep=6pt, halign=c, valign=m},
  row{1} = {bg=StageBlue!15, fg=black, font=\bfseries, valign=m},
  row{odd} = {bg=white},
  row{even} = {bg=LightGray},
}
\Hdr{Phase} &
\Hdr{Required Tier (from Stage 6)} &
\Hdr{What Changes / Becomes Mandatory} &
\Hdr{Notes} \\
Phase 00 &
T0 &
None &
Optional-layer only. Coffee/tea (Tier 0) may exist, but no novelty near protected windows. If prescription stimulant use exists, OTC stimulant additions are clinician-led only. \\
Phase 01 &
T0 &
None &
Optional-layer only. Stable variables only; no new stimulants introduced in Deload+Test windows. \\
Phase 02 &
T0 &
None &
Optional-layer only. Do not use stimulants to compensate for increased \textbf{STR} stress or sleep debt; downshift dose instead. \\
Phase 03 &
T0 &
None &
Optional-layer only. Ruck introduction phases increase fatigue cost; respond with recovery controls, not stimulant escalation. \\
Phase 04 &
T0 &
None &
Optional-layer only. Aquatic relevance does not change stimulant posture; maintain stability and documentation. \\
Phase 05 &
T0 &
None &
Optional-layer only. Longer endurance windows increase the cost of GI/sleep disruption; omit is preferred if stability is not proven. \\
Phase 06 &
T0 &
None &
Optional-layer only. Higher consequence windows require low variance; any change near gates is discouraged and logged as a controlled variable if it occurs. \\
Phase 07 &
T0 &
None &
Optional-layer only. Hybrid density increases risk of stimulant dependence; enforce fundamentals-first and sleep protection. \\
Phase 08 &
T0 &
None &
Optional-layer only. Durability/load tolerance phases treat sleep and recovery as primary controls; stimulants remain optional and conservative. \\
Phase 09 &
T0 &
None &
Optional-layer only. High \textbf{CAL} volume phases should not introduce stimulant novelty near protected windows. \\
Phase 10 &
T0 &
None &
Optional-layer only. If any testing/clearance implications exist, contamination risk posture becomes more important (COA preference; avoid blends). \\
Phase 11 &
T0 &
None &
Optional-layer only. Late phases reward low variance and stable sleep; omit is preferred if any instability exists. \\
Phase 12 &
T0 &
None &
Optional-layer only. Verification blocks require stable variables; no changes near protected windows. \\
Phase 13 &
T0 &
None &
Optional-layer only. Consolidation demands repeatability; stimulants do not create gate credit and do not justify advancement claims. \\
\end{longtblr}

\endgroup

Notes:

\begin{itemize}
\item Phase Integration Map displays Stage 6 required tier values as-is (T0 throughout). Any downstream artifact implying 7.05 is required for deployability is a mismatch and requires Stage 8 review.
\end{itemize}

\subsection*{7.05.9 Risks, Contraindications, and Red Flags}
\phantomsection
\addcontentsline{toc}{subsection}{7.05.9 Risks, Contraindications, and Red Flags}

\begin{itemize}
\item Anxiety/panic escalation risk (hard stop):
  \begin{itemize}
  \item panic symptoms, severe anxiety spike, agitation, or dissociation triggers: stop, log, do not test that day; clinician referral posture when severe or recurrent.
  \end{itemize}
\item Cardiovascular sensitivity risk (hard stop when severe):
  \begin{itemize}
  \item palpitations with dizziness, chest pressure/pain, near-syncope, or disproportionate shortness of breath triggers \textbf{STOP} \textbf{NOW} posture and escalation per governance.
  \end{itemize}
\item Sleep displacement risk:
  \begin{itemize}
  \item stimulants can cause delayed sleep onset and sleep fragmentation; this degrades recovery and invalidates comparability.
  \item if sleep collapses after introduction, stop the new variable and revert to baseline.
  \end{itemize}
\item Medication interaction risk (referral-only posture):
  \begin{itemize}
  \item if on prescription stimulants or blood pressure medication, OTC stimulant use is clinician-led only; default omit unless cleared.
  \item if on psychiatric medications, treat activating OTC products as higher risk; clinician discussion when uncertain; omit if ambiguous.
  \end{itemize}
\item Contamination and banned-substance risk posture:
  \begin{itemize}
  \item proprietary blends and undisclosed stimulants increase contamination risk and cannot be defended in an audit posture.
  \item control: avoid blends; prefer transparent labeling and third-party testing/COA when available.
  \end{itemize}
\item Prohibited product risk:
  \begin{itemize}
  \item fat burners, DMAA/DMHA-like, ephedra-like, “research” nootropics, and undisclosed stimulants are prohibited; recommending them is non-deployable and requires Stage 8 review.
  \end{itemize}
\end{itemize}

\subsection*{7.05.10 Compliance Checklist — Pass/Fail}
\phantomsection
\addcontentsline{toc}{subsection}{7.05.10 Compliance Checklist — Pass/Fail}

\begin{itemize}
\item \textbf{PASS}/\textbf{FAIL}: OTC-only posture maintained; no prescription guidance, no hormone guidance, no illegal/gray-market substances.
\item \textbf{PASS}/\textbf{FAIL}: No dosing ranges, cycling protocols, stack prescriptions, or test-day dosing instructions appear.
\item \textbf{PASS}/\textbf{FAIL}: Tier 0 includes coffee and tea (household) and Tier 1 includes single-ingredient caffeine formats; no proprietary blends promoted.
\item \textbf{PASS}/\textbf{FAIL}: Prohibited product classes are explicitly listed (proprietary stimulant blends, fat burners, DMAA/DMHA-like, ephedra-like, undisclosed stimulants).
\item \textbf{PASS}/\textbf{FAIL}: Hard stop triggers include panic/anxiety spike with “stop, log, do not test that day” posture.
\item \textbf{PASS}/\textbf{FAIL}: Interaction caution posture is explicit: prescription stimulant or BP medication context is clinician-led only; no medical advice given.
\item \textbf{PASS}/\textbf{FAIL}: Change-control posture is explicit: one item at a time; no novelty near Deload+Test/gate windows; logging required.
\item \textbf{PASS}/\textbf{FAIL}: Phase Integration Map reflects Stage 6 required tier values as-is and does not frame 7.05 as required for deployability.
\end{itemize}

\end{document}